%%%%%%%%%%%%%%%%%%%%%%%%%%%%%%%%%%%%%%%%%%%%%%%%%%%%%%%%%%%%%%%%%%%%%%
% Lớp 11 | Xác suất | Biến cố hợp và quy tắc cộng xác suất
% Dạng 1D9?2-3: Quy tắc cộng cho hai biến cố xung khắc
% Số câu: 33  (sinh tự động bởi tach_cau.py; từ nay thầy sửa trực tiếp file này)
%%%%%%%%%%%%%%%%%%%%%%%%%%%%%%%%%%%%%%%%%%%%%%%%%%%%%%%%%%%%%%%%%%%%%%

%%%%%%%%%%%%%%%% TRẮC NGHIỆM 4 PHƯƠNG ÁN %%%%%%%%%%%%%%%%

% nguồn: Ok_TracNghiemLop11_Lan1
\begin{ex}%[1D9N2-3]%[ma:e8fb1491]
	Do con súc sắc không cân đối nên mỗi khi gieo  con súc sắc đó thì xác suất xuất hiện mặt $2$ chấm và mặt $5$ chấm lần lượt là $0{,}35$ và $0{,}25$. Xác suất xuất hiện mặt $2$ chấm hoặc mặt $5$ chấm là
	 \choice
		 {\True $0{,}6$}
		 {$0{,}0875$}
		{$0{,}4875$} 
		{$0{,9125}$} 
	\loigiai{ 
		Gọi $A$ là biến cố xuất hiện  mặt $2$ chấm khi gieo con súc sắc.\\
		Gọi $B$ là biến cố xuất hiện mặt $5$ chấm  khi gieo con súc sắc.\\
		Do đó, $A\cup B$ là biến cố xuất hiện mặt $2$ chấm  hoặc mặt $5$ chấm.\\
		Ta có $\mathrm{P}(A)=0{,}35$  và $\mathrm{P}(B)=0{,}25$.\\
		 Do $A$ và $B$ là hai biến cố xung khắc nhau nên
		 \[\mathrm{P}(A\cup B)=\mathrm{P}(A)+\mathrm{P}(B)=0{,}35+0{,}25=0{,}6.\]
	} 
\end{ex}

% nguồn: Ok_TracNghiemLop11_Lan1
\begin{ex}%[1D9N2-3]%[ma:24aa6226]
	Cho $A$ và $B$ là hai biến cố độc lập. Biết $\mathrm{P}(A) = 0{,}4$ và  $\mathrm{P}(B) = 0{,}5$. Xác suất của biến cố $A \cup B$ là
	\choice 
		{\True $0{,}9$} 
		{$0{,}7$} 
		{$0{,}5$} 
		{$0{,}2$} 
	 \loigiai{
		Do $A$ và $B$ là hai biến cố độc lập nên nó là hai  biến cố xung khắc.\\
		Suy ra $\mathrm{P}\left(A \cup B\right)=\mathrm{P}(A)+\mathrm{P}(B)=0{,}9$. 
	} 
\end{ex}

% nguồn: Ok_TracNghiemLop11_Lan1
\begin{ex}%[1D9H2-3]%[ma:78378eb4]
	Một hộp chứa $6$ quả cầu trắng và $7$ quả cầu đen. Lấy ngẫu nhiên đồng thời bốn quả trong số đó. Xác suất để có ít nhất một quả  cầu trắng bằng
	 \choice
		 {$\dfrac{7}{143}$}
		{$\dfrac{120}{143}$} 
		 {$\dfrac{13}{24}$}
		{\True $\dfrac{136}{143}$} 
	 \loigiai{
		Số phần tử của không gian mẫu là  $n(\Omega)=\mathrm{C}_{13}^4$.\\
		Gọi $A$ là biến cố lấy được ít nhất một quả  cầu trắng.\\
		Biến cố đối của $A$ là biến cố  $B$ \lq\lq không lấy được quả cầu trắng nào\rq\rq.\\
		 Số phần tử của biến cố $B$ là $n(B)=\mathrm{C}_7^4$.\\
		Xác  suất $\mathrm{P}(B)=\dfrac{n(B)}{n(\Omega)}=\dfrac{\mathrm{C}_7^4}{\mathrm{C}_{13}^4}=\dfrac{7}{143}$.\\
		Suy ra  xác suất $\mathrm{P}(A)=1-\mathrm{P}(B)=\dfrac{136}{143}$.
	} 
\end{ex}

% nguồn: Ok_TracNghiemLop11_Lan1
\begin{ex}%[1D9H2-3]%[ma:5213456b]
	Lớp có $40$ học sinh gồm $16$ nam và $24$ nữ. Thầy giáo chủ nhiệm cần chọn ngẫu nhiên trong lớp một nhóm gồm $3$ học sinh để tham gia văn nghệ kỉ niệm ngày thành lập  trường. Tính xác suất nhóm thầy chọn được có ít nhất $2$ học sinh nữ.
	 \choice
		 {$\dfrac{193}{247}$}
		{$\dfrac{552}{1235}$} 
		{$\dfrac{253}{1235}$} 
		{\True  $\dfrac{161}{247}$}
	\loigiai{ 
		 Ta có $n(\Omega)=\mathrm{C}_{40}^3$.\\
		Xét biến cố $A\colon$ \lq\lq $3$ học sinh để tham gia văn nghệ có ít nhất  $2$ học sinh nữ\rq\rq.\\
		 Suy ra $n(A)=\mathrm{C}_{16}^1\cdot\mathrm{C}_{24}^2+\mathrm{C}_{24}^3$.\\
	Xác suất  $\mathrm{P}(A)=\dfrac{n(A)}{n(\Omega)}=\dfrac{161}{247}$.}
\end{ex}

% nguồn: Ok_TracNghiemLop11_Lan1
\begin{ex}%[1D9H2-3]%[ma:0dd37d26]
	Hai bạn An và Bình độc lập với nhau tham gia một cuộc thi. Xác suất để bạn An và bạn Bình đạt giải tương ứng là $0{,}8$ và $0{,}6$. Xác suất để có ít nhất một bạn đạt giải  là
	 \choice
		{$0{,}94$} 
		 {$0{,}924$}
		 {\True $0{,}92$}
		{$0{,}93$} 
	\loigiai{ 
		Xác suất để bạn An và bạn Bình đạt giải tương ứng  là $0{,}8$ và $0{,}6$.\\
		Suy ra xác suất để bạn An và bạn Bình không đạt giải tương ứng là $0{,}2$ và $0{,}4$.\\ 
		 Xác suất để có không bạn nào đạt giải là $0{,}2\cdot 0{,}4=0,08$.\\
		Vậy xác suất để có ít nhất một bạn đạt  giải là $1-0{,}08=0{,}92$.
	 }
\end{ex}

% nguồn: Ok_TracNghiemLop11_Lan1
\begin{ex}%[1D9V2-3]%[ma:a6b82adf]
	Trong một lớp học gồm $15$ học sinh  nam và $10$ học sinh nữ. Giáo viên gọi ngẫu nhiên $4$ học sinh lên giải bài tập. Xác suất để $4$ học sinh được gọi đó có cả nam và nữ là
	\choice 
		{$\dfrac{219}{233}$} 
		 {$\dfrac{442}{506}$}
		{\True $\dfrac{443}{506}$} 
		{$\dfrac{219}{323}$} 
	\loigiai{ 
		 Số cách chọn $4$ học sinh từ $25$ học sinh là $\mathrm{C}_{25}^4$.\\
		Số cách chọn $4$ học sinh trong đó chỉ có  nam hoặc nữ là $\mathrm{C}_{15}^4+\mathrm{C}_{10}^4$.\\
	Xác suất để $4$ học sinh  được gọi có cả nam và nữ là $1-\dfrac{\mathrm{C}_{15}^4+\mathrm{C}_{10}^4}{\mathrm{C}_{25}^4}=\dfrac{443}{506}$.}
\end{ex}


%%%%%%%%%%%%%%%% ĐÚNG/SAI %%%%%%%%%%%%%%%%

% nguồn: Ok_dungsai_lop11_Lan1
\begin{ex}%[1D9H2-3]%[ma:24ca84fc]
 Lớp $10A1$ có $35$ học sinh gồm $20$ học sinh nữ và $15$ học sinh nam. Nhà trường cần chọn $3$ học  sinh của lớp đó tham gia trồng cây. Các khẳng định sau đây, khẳng định nào là đúng, khẳng định nào là sai?
  \choiceTF
  {\True Xác suất để $3$ học sinh được chọn có đúng  một học sinh nữ bằng $\dfrac{60}{187}$}
  {\True Xác suất để $3$ học sinh được chọn có đúng hai học sinh nữ bằng  $\dfrac{570}{1309}$}
  {Xác suất để $3$ học sinh được  chọn có ít nhất một học sinh nữ bằng $\dfrac{90}{119}$}
  {Xác suất để $3$ học sinh được chọn có nhiều nhất hai học sinh nữ  bằng $\dfrac{114}{187}$}
 \loigiai{ 
  Không gian  mẫu $n\left(\Omega \right)=\mathrm{C}_{35}^3$.
  \begin{itemchoice}
   \itemch $A\colon$ \lq\lq $3$ học sinh được chọn có đúng một học sinh nữ\rq\rq. Ta có $n\left(A \right)=20\cdot\mathrm{C}_{15}^2$.\\ 
   Xác suất của biến  cố $A$ là $\mathrm{P}\left(A \right)=\dfrac{n\left(A \right)}{n\left(\Omega \right)}=\dfrac{20\cdot\mathrm{C}_{15}^2}{\mathrm{C}_{35}^3}=\dfrac{60}{187}$.\\
   \itemch $B\colon$  \lq\lq $3$ học sinh được chọn có đúng hai học sinh nữ\rq\rq. Ta có $n\left(B \right)=15\cdot\mathrm{C}_{20}^2$.\\
   Xác suất của biến cố $B$ là $\mathrm{P}\left(B \right)=\dfrac{n\left(B \right)}{n\left(\Omega  \right)}=\dfrac{15\cdot\mathrm{C}_{20}^2}{\mathrm{C}_{35}^3}=\dfrac{570}{1309}$.\\
   \itemch $C\colon$ \lq\lq $3$ học sinh được chọn có ít nhất một học sinh  nữ\rq\rq. Ta có $n\left(C \right)=20\cdot\mathrm{C}_{15}^2+15\cdot\mathrm{C}_{20}^2+\mathrm{C}_{20}^3$.\\
   Xác suất của biến cố $C$ là $\mathrm{P}\left(C \right)=\dfrac{n\left(C \right)}{n\left(\Omega \right)}=\dfrac{20\cdot\mathrm{C}_{15}^2+15\cdot\mathrm{C}_{20}^2+\mathrm{C}_{20}^3}{\mathrm{C}_{35}^3}=\dfrac{174}{187}$. 
   \itemch $D\colon$ \lq\lq $3$ học sinh được chọn có nhiều nhất hai học sinh nữ\rq\rq. Ta có $n\left(D  \right)=\mathrm{C}_{15}^3+20\cdot\mathrm{C}_{15}^2+15\cdot\mathrm{C}_{20}^2$.\\
   Xác suất của biến cố $D$ là $\mathrm{P}\left(D \right)=\dfrac{n\left(D \right)}{n\left(\Omega \right)}=\dfrac{\mathrm{C}_{15}^3+20\cdot\mathrm{C}_{15}^2+15\cdot\mathrm{C}_{20}^2}{\mathrm{C}_{35}^3}=\dfrac{1081}{1309}$. 
  \end{itemchoice}
  }
\end{ex}

% nguồn: Ok_dungsai_lop11_Lan1
\begin{ex}%[1D9H2-3]%[ma:70986349]
 Một hộp đựng $4$ viên bi màu xanh, $3$ viên bi màu đỏ và $2$ viên bi màu vàng. Chọn ngẫu nhiên $2$ viên bi từ hộp trên. Gọi $A$ là biến cố ``Chọn được $2$ viên bi màu xanh'', $B$ là biến cố ``Chọn được $2$ viên bi màu đỏ'', $C$ là biến cố ``Chọn được 2 viên bi  màu vàng''. Xét tính Đúng - Sai của các mệnh đề sau:
  \choiceTF{$\mathrm{P}(A)=\dfrac{1}{7}$}
  {$\mathrm{P}(B)=\dfrac{1}{8}$} 
  {\True $\mathrm{P}(C)=\dfrac{1}{36}$} 
  {\True Xác suất  để chọn được hai viên bi cùng màu bằng $\dfrac{5}{18}$}
  \loigiai{
  \begin{itemchoice}
    \itemch $\mathrm{P}(A)=\dfrac{\mathrm{C}^2_4}{\mathrm{C}^2_9}=\dfrac{1}{6}$.
   \itemch  $\mathrm{P}(B)=\dfrac{\mathrm{C}^2_3}{\mathrm{C}^2_9}=\dfrac{1}{12}$.
   \itemch  $\mathrm{P}(C)=\dfrac{\mathrm{C}^2_2}{\mathrm{C}^2_9}=\dfrac{1}{36}$.
   \itemch Gọi $X$ là biến cố ``Chọn được  hai bi cùng màu''.\\
   Ta có $X = A \cup B \cup C$ và $A$, $B$,  $C$ là các biến cố đôi một xung khắc.\\
    Do đó, $\mathrm{P}(X)=\mathrm{P}(A)+\mathrm{P}(B)+\mathrm{P}(C)=\dfrac{1}{6}+\dfrac{1}{12}+ \dfrac{1}{36}=\dfrac{5}{18}$.
  \end{itemchoice}
  }
\end{ex}

% nguồn: Ok_dungsai_lop11_Lan1
\begin{ex}%[1D9H2-3]%[ma:8ca4883a]
 Một hộp đựng $15$ tấm thẻ cùng loại được đánh số từ $1$ đến $15$. Rút ngẫu nhiên một tấm  thẻ trong hộp. Xét các biến cố sau\\
 \indent  \indent $A$ là biến cố \lq\lq Số ghi trên tấm thẻ là số chẵn\rq\rq; \\
 \indent \indent $B$ là biến cố \lq\lq Số ghi trên tấm thẻ là số chính phương\rq\rq;  \\
 \indent \indent $C$ là biến cố \lq\lq Số ghi trên thẻ là số chẵn hoặc số chính  phương\rq\rq.\\
 Các khẳng định sau  đúng hay sai?
  \choiceTF
   {\True $n \left(\Omega\right)=15$}
   {\True $C=A\cup B$}
   {$\mathrm{P}\left(A\right) = \dfrac{8}{15}$}
  {\True  $\mathrm{P}\left(C\right) = \dfrac{2}{3}$}
  \loigiai{
  \begin{itemchoice}
    \itemch Không gian mẫu $\Omega=\left\{1;2;3;4;5;\cdots;14;15\right\} \Rightarrow n \left(\Omega\right)=15$.
   \itemch $A=\left\{2;4;6;8;10;12;14\right\}$; $B=\left\{1;4;9\right\}$;  $C=\left\{1;2;4;6;8;9;10;12;14\right\}$.\\
   Vậy $C=A\cup B$. 
   \itemch Ta có $n(A)=  7$; $n \left(\Omega\right)=15 \Rightarrow \mathrm{P}\left(A\right) = \dfrac{7}{15}$.
   \itemch Ta có $n(A)= 8$; $n \left(\Omega\right)=15 \Rightarrow  \mathrm{P}\left(B\right) = \dfrac{8}{15}$.\\
   Do đó $\mathrm{P}\left(C\right) = \mathrm{P}\left(A\right)+\mathrm{P}\left(B\right) = \dfrac{2}{3}$. 
  \end{itemchoice}
  }
\end{ex}

% nguồn: Ok_dungsai_lop11_Lan1
\begin{ex}%[1D9H2-3]%[ma:db916525]
 Một đội tình nguyện gồm $6$ học sinh khối $10$, $5$ học sinh khối $11$ và $4$ học sinh khối $12$. Chọn ngẫu nhiên $5$ học sinh trong đội. Các khẳng  định sau đây, khẳng định nào đúng, khẳng định nào sai?
  \choiceTF
   {Số phần tử không gian mẫu là $15$}
   {\True Số phần tử của biến cố \lq\lq cả $5$ học sinh được chọn học cùng một khối\rq\rq~ là $\mathrm{C}_6^5+\mathrm{C}_5^5$}
  {Xác suất của biến cố \lq\lq cả $5$ học sinh được chọn học cùng  một khối\rq\rq~ là $\dfrac{5}{3003}$}
  {\True Xác suất của biến cố  \lq\lq chọn được $5$ học sinh đủ ba khối\rq\rq~ là $\dfrac{310}{429}$}
 \loigiai{ 
  \begin{itemchoice}
    \itemch Số phần tử không gian mẫu là $n(\Omega)=\mathrm{C}_{15}^5=3003$.
    \itemch Gọi $A$ là biến cố \lq\lq cả $5$ học sinh được chọn học cùng một khối\rq\rq.\\
   Vì cả $5$ học sinh được chọn học cùng một khối nên chỉ có thể chọn $5$ học sinh khối $10$ hoặc $5$ học sinh  khối $11$.\\
   Suy ra  $n(A)=\mathrm{C}_6^5+\mathrm{C}_5^5$.
   \itemch Xác suất của biến cố \lq\lq cả $5$ học  sinh được chọn học cùng một khối\rq\rq~ là
   $$\mathrm{P}(A)=\dfrac{n(A)}{n(\Omega)}=\dfrac{\mathrm{C}_6^5+\mathrm{C}_5^5}{\mathrm{C}_{15}^5}=\dfrac{1}{429}.$$ 
   \itemch Gọi $B$ là biến cố \lq\lq chọn được $5$ học  sinh đủ ba khối\rq\rq.\\
   Khi  đó $n(B)=\mathrm{C}_6^1\mathrm{C}_5^1\mathrm{C}_4^3+\mathrm{C}_6^1\mathrm{C}_5^3\mathrm{C}_4^1+\mathrm{C}_6^1\mathrm{C}_5^2\mathrm{C}_4^2+\mathrm{C}_6^2\mathrm{C}_5^1\mathrm{C}_4^2+\mathrm{C}_6^2\mathrm{C}_5^2\mathrm{C}_4^1+\mathrm{C}_6^3\mathrm{C}_5^1\mathrm{C}_4^1=2170$.\\
    Vậy $\mathrm{P}(B)=\dfrac{n(B)}{n(\Omega)}=\dfrac{2170}{\mathrm{C}_{15}^5}=\dfrac{310}{429}$.
  \end{itemchoice}
  }
\end{ex}

% nguồn: Ok_dungsai_lop11_Lan1
\begin{ex}%[1D9H2-3]%[ma:c308c464]
 Một hộp chứa $5$ quả bóng xanh, $6$ quả bóng đỏ và $2$  quả bóng vàng có cùng kích thước và khối lượng. Chọn ra ngẫu nhiên từ hộp $3$ quả bóng. Khẳng định nào sau đây đúng, khẳng định nào sau đây sai?
 \choiceTF 
  {\True Số phần tử không gian mẫu là  $286$ phần tử}
   {Số phần tử của biến cố \lq\lq Chọn được $3$ quả bóng khác màu\rq\rq~ là $13$}
  {Xác suất của biến cố \lq\lq Cả  $3$ quả bóng lấy ra đều có cùng màu\rq\rq~ là $\dfrac{113}{143}$}
  {\True Xác suất của biến cố  \lq\lq Có đúng $2$ quả bóng xanh trong $3$ quả bóng lấy ra\rq\rq~ là $\dfrac{45}{143}$}
 \loigiai{ 
  \begin{itemchoice}
   \itemch Chọn ngẫu nhiên từ hộp $3$ quả bóng trong tổng số $13$  quả bóng có $\mathrm{C}_{13}^3=286$ cách.\\
   Vậy $n(\Omega)=286$ phần  tử.
   \itemch Gọi $T\colon$ \lq\lq Chọn được $3$ quả bóng khác màu\rq\rq.\\ 
   Vậy  chọn $1$ quả bóng xanh, $1$ quả bóng đỏ và $1$ quả bóng vàng có $5\cdot 6\cdot 2=60$ cách.
   \itemch Gọi $A$ là biến cố  \lq\lq Cả $3$ quả bóng lấy ra đều có cùng màu xanh\rq\rq.\\
   Do đó  $n(A)=\mathrm{C}_5^3=10 \Rightarrow \mathrm{P}(A)=\dfrac{n(A)}{n(\Omega)}=\dfrac{10}{286}=\dfrac{5}{143}$.\\
   Gọi $B$ là biến cố \lq\lq Cả  $3$ quả bóng lấy ra đều có cùng màu đỏ\rq\rq.\\
    Do đó $n(B)=\mathrm{C}_6^3=20 \Rightarrow \mathrm{P}(B)=\dfrac{n(B)}{n(\Omega)}=\dfrac{20}{286}=\dfrac{10}{143}$.\\
   Gọi $C$ là biến cố \lq\lq Cả $3$ quả bóng lấy ra đều có  cùng màu vàng\rq\rq.\\
   Do  đó $n(C)=0 \Rightarrow \mathrm{P}(C)=0$.\\
   Ta có $A\cup B\cup C$ là biến cố \lq\lq Cả $3$ quả bóng lấy ra đều có cùng màu\rq\rq.\\ 
   Suy  ra $\mathrm{P}(A\cup B\cup C)=\mathrm{P}(A)+\mathrm{P}(B)+\mathrm{P}(C)=\dfrac{15}{243}$.
   \itemch Gọi $D$ là biến cố \lq\lq Có đúng $2$ quả bóng  xanh trong $3$ quả bóng lấy ra\rq\rq.\\
   Do đó $A\cup D$ là biến cố \lq\lq Có ít nhất $2$ quả bóng xanh trong $3$ quả bóng lấy ra\rq\rq.\\ 
   Chọn ngẫu nhiên từ hộp $2$ quả bóng trong tổng số $5$ quả bóng xanh có $\mathrm{C}_5^2=10$  cách.\\
   Chọn  ngẫu nhiên từ hộp $1$ quả bóng trong tổng số $8$ quả bóng đỏ hoặc vàng có $\mathrm{C}_8^1=8$ cách.\\
    Suy ra $n(D)=10\cdot 8=80 \Rightarrow \mathrm{P}(D)=\dfrac{n(D)}{n(\Omega)}=\dfrac{80}{286}=\dfrac{40}{143}$.\\
   Vậy $\mathrm{P}(A\cup D)=\mathrm{P}(A)+\mathrm{P}(D)=\dfrac{45}{143}$. 
  \end{itemchoice}
  }
\end{ex}

% nguồn: Ok_dungsai_lop11_Lan1
\begin{ex}%[1D9H2-3]%[ma:c91a090d]
 Một hộp có $12$ chiếc thẻ cùng loại, mỗi thẻ được ghi một trong các số $1,2,3,  \ldots, 12$; hai thẻ khác nhau thì ghi hai số khác nhau. Rút ngẫu nhiên $1$ chiếc thẻ trong hộp. Xét
 \begin{itemize}
  \item Biến cố $A$ \lq\lq Số xuất  hiện trên thẻ được rút ra là số chia hết cho $3$\rq\rq.
  \item Biến cố $B$ \lq\lq Số xuất hiện trên thẻ được rút ra là  số chia hết cho $5$\rq\rq.
 \end{itemize}
 \choiceTF 
  {\True Số phần tử của  không gian mẫu $n(\Omega)=12$}
  {\True Xác xuất của  biến cố $A$ là $\mathrm{P}(A)=\dfrac{1}{3}$}
  {\True Xác xuất của biến  cố $B$ là $\mathrm{P}(B)=\dfrac{1}{6}$}
  {Xác xuất của biến cố hợp là $\mathrm{P}(A \cup  B)=\dfrac{1}{18}$}
 \loigiai{ 
  \begin{itemchoice}
   \itemch Không gian mẫu của phép thử trên  có $12$ phần tử, tức là $n(\Omega)=12$.
   \itemch Số các kết quả thuận  lợi cho các biến cố $A$ là $n(A)=4$.\\
   Suy ra $\mathrm{P}(A)=\dfrac{n(A)}{n(\Omega)}=\dfrac{4}{12}=\dfrac{1}{3}$. 
   \itemch Số các kết quả thuận lợi cho các biến cố $B$  à $n(B)=2$.\\
   Suy ra $\mathrm{P}(B)=\dfrac{n(B)}{n(\Omega)}=\dfrac{2}{12}=\dfrac{1}{6}$. 
   \itemch Trong các  số $1,2,3, \ldots, 12$, không có số nào chia hết cho cả $3$ và $5$.\\
   Vì thế $A,  B$ là hai biến cố xung khắc.\\
   Suy ra $\mathrm{P}(A \cup B)=\mathrm{P}(A)+\mathrm{P}(B)=\dfrac{1}{3}+\dfrac{1}{6}=\dfrac{1}{2}$. 
  \end{itemchoice}
 } 
\end{ex}

% nguồn: Ok_dungsai_lop11_Lan1
\begin{ex}%[1D9H2-3]%[ma:7ed8c508]
 Một vận động viên thi bắn súng. Biết rằng xác suất để vận động viên đó bắn trúng vòng $10$ là  $0{,}25$; bắn trúng vòng $9$ là $0{,}3$; bắn trúng vòng $8$ là $0{,}4$. Nếu bắn trúng vòng $k$ thì được $k$ điểm. Vận động viên thực hiện bắn hai lần, hai lần bắn độc lập với nhau. Các khẳng định sau đây, khẳng định nào đúng, khẳng định nào sai?
  \choiceTF
  {Xác suất để cả hai lần bắn của vận động viên đó đều trúng vòng $9$  là $0{,}6$}
  {Xác suất để vận  động viên đó đạt $16$ điểm là $0{,}8$}
  {\True Xác suất để vận động viên  đó đạt $17$ điểm là $0{,}24$}
  {Xác suất để vận động viên đó  có số điểm lớn hơn $17$ là $0{,}3025$}
 \loigiai{ 
  $A\colon$ \lq\lq vận động viên đó bắn trúng vòng  $10$\rq\rq~ $\Rightarrow \mathrm{P}(A)=0{,}25$.\\
  $B\colon$ \lq\lq vận động viên đó bắn trúng vòng  $9$\rq\rq~ $\Rightarrow \mathrm{P}(B)=0{,}3$.\\
  $C\colon$ \lq\lq vận động viên đó bắn trúng vòng $8$\rq\rq~  $\Rightarrow \mathrm{P}(C)=0{,}4$.
  \begin{itemchoice}
    \itemch Ta có xác suất để cả hai lần bắn đều trúng vòng $9$ là $\mathrm{P}(BB)=0{,}3\cdot 0{,}3=0,09$.
   \itemch Ta có vận động viên đó đạt $16$ điểm thì cả hai lần bắn đều trúng vòng  $8$.\\
   Xác suất là $\mathrm{P}(CC)=0{,}4\cdot  0{,}4=0,16$.\\
   \itemch Ta có vận động viên đó đạt $17$ điểm thì lần $1$ trúng bắn trúng vòng $9$ và lần $2$ bắn trúng vòng $8$ hoặc lần $1$ trúng bắn trúng vòng $8$ và lần  $2$ bắn trúng vòng $9$.\\
   Xác suất là  $\mathrm{P}(BC\cup CB)=0{,}3\cdot 0{,}4+0{,}4\cdot 0{,}3=0{,}24$.
   \itemch Ta có vận động viên đó có số điểm lớn hơn $17$ thì xác suất là 
   $$\mathrm{P}(AA\cup AB\cup  BA\cup BB\cup AC\cup CA)=0{,}25\cdot 0{,}25+0{,}25\cdot 0{,}3+0{,}3\cdot 0{,}25+0{,}3\cdot 0{,}3+0{,}25\cdot 0{,}4+0{,}4\cdot 0{,}25=0{,}5025.$$
  \end{itemchoice}
 } 
\end{ex}

% nguồn: Ok_dungsai_lop11_Lan1
\begin{ex}%[1D9H2-3]%[ma:50ad473e]
 Một đội tình nguyện gồm $9$ học sinh khối $10$ và $7$ học sinh  khối $11$. Chọn ra ngẫu nhiên $3$ người trong đội. Xét
 \begin{itemize}
  \item $A$  là biến cố \lq\lq Cả $3$ học sinh được chọn đều thuộc khối $10$ \rq\rq
  \item $B$ là biến cố \lq\lq Cả $3$ học sinh  được chọn đều thuộc khối $11$\rq\rq.
 \end{itemize}
  \choiceTF
  {\True  Số phần tử của không gian mẫu $n(\Omega)=\mathrm{C}_{16}^3$}
   {Xác xuất của biến cố $A$ là $\mathrm{P}(A)=\dfrac{\mathrm{C}_7^3}{\mathrm{C}_{16}^3}$}
  {Xác xuất của biến cố $B$ là  $\mathrm{P}(B)=\dfrac{\mathrm{C}_9^3}{C_{16}^3}$}
  {\True Xác suất của biến cố \lq\lq  Cả $3$ người được chọn học cùng một khối\rq\rq\, là $\dfrac{17}{80}$}
  \loigiai{
  \begin{itemchoice}
    \itemch Ta có $n(\Omega)=\mathrm{C}_{16}^3$.
   \itemch Gọi $A$ là biến cố \lq\lq Cả $3$ học sinh được chọn đều thuộc  khối $10$\rq\rq.\\
   Khi  đó, $\mathrm{P}(A)=\dfrac{\mathrm{C}_9^3}{\mathrm{C}_{16}^3}$.
   \itemch $B$ là biến cố \lq\lq Cả $3$  học sinh được chọn đều thuộc khối $11$\rq\rq.\\
    Khi đó, $\mathrm{P}(B)=\dfrac{\mathrm{C}_7^3}{C_{16}^3}$.
   \itemch  $A \cup B$ là biến cố \lq\lq Cả $3$ người được chọn học cùng một khối\rq\rq.\\
   Do $A$ và $B$ là hai biến cố xung khắc nên $\mathrm{P}(A  \cup B)=\mathrm{P}(A)+\mathrm{P}(B)$.\\
    Vậy, $\mathrm{P}(A \cup B)=\dfrac{\mathrm{C}_9^3+\mathrm{C}_7^3}{\mathrm{C}_{16}^3}=\dfrac{17}{80}$.
  \end{itemchoice}
 } 
\end{ex}

% nguồn: Ok_dungsai_lop11_Lan1
\begin{ex}%[1D9H2-3]%[ma:f9485c6d]
 Cho tập $E=\{1,2,3,4,5\}$. Viết ngẫu nhiên lên bảng hai số tự nhiên, mỗi số gồm $3$ chữ số đôi một khác nhau từ tập $E$. Gọi $A$ là biến cố \lq\lq Số viết trước có chữ số $5$ và số viết sau không có chữ số $5$\rq\rq,  $B$ là biến cố \lq\lq Số viết trước không có chữ số $5$ và số viết sau có chữ số $5$\rq\rq. Trong các khẳng định sau, khẳng định nào đúng, khẳng định nào sai?
  \choiceTF
  {\True $A$ và  $B$ là hai biến cố xung khắc}
  {\True $A\cup B$ là biến cố \lq\lq Trong hai số đó có đúng một số có chữ số $5$\rq\rq} 
   {$\mathrm{P}(A)=\dfrac{865}{3600}$}
  {\True $\mathrm{P}(A\cup B)=\dfrac{12}{25}$} 
  \loigiai{
  Gọi $S$ là tập hợp các số tự nhiên gồm $3$ chữ số phân biệt lập từ tập $E$ thì số phần tử của $S$ là  $n(S)=\mathrm{A}_5^3=60$.\\
  Gọi $F$ là tập  hợp các số tự nhiên gồm $3$ chữ số phân biệt lập từ tập $E$ sao cho trong số đó có mặt chữ số $5$.\\
  Do  đó, số phần tử của $F$ là $n(F)=3\cdot\mathrm{A}_4^2=36$.\\
  Không gian mẫu $\Omega$ của phép thử trên có số  phần tử là $n\left(\Omega \right)=60\cdot 60=3600$.
  \begin{itemchoice}
   \itemch Hai biến cố $A$ và $B$  không cùng xảy ra trong một lần thử, do đó $A$ và $B$ là hai biến cố xung khắc.
   \itemch $A\cup B$ là biến cố  \lq\lq Số viết trước có chữ số $5$, số viết sau không có chữ số $5$ hoặc số viết trước không có chữ số $5$, số viết sau có chữ số $5$\rq\rq.\\
   Hay $A\cup B$ là biến cố \lq\lq Trong hai số đó có đúng một số có chữ số  $5$\rq\rq.
   \itemch  Ta có $n\left(A\right)=24\cdot 36=864 \Rightarrow \mathrm{P}(A)=\dfrac{n\left(A\right)}{n\left(\Omega \right)}=\dfrac{864}{3600}$.
   \itemch  Ta có $n\left(B\right)=36\cdot 24=864 \Rightarrow \mathrm{P}(B)=\dfrac{n\left(B\right)}{n\left(\Omega \right)}=\dfrac{864}{3600}$.\\
   Vậy  $\mathrm{P}(A\cup B)=\dfrac{864}{3600}+\dfrac{864}{3600}=\dfrac{12}{25}$.
  \end{itemchoice}
  }
\end{ex}

% nguồn: Ok_dungsai_lop11_Lan1
\begin{ex}%[1D9H2-3]%[ma:83563a70]
 Gieo con xúc xắc  cân đối và đồng chất $1$ lần. Khẳng định nào sau đây đúng, khẳng định nào sau đây sai?
  \choiceTF
  {\True Số chấm xuất hiện là số lẻ thì xác xuất của nó  bằng $\dfrac{1}{2}$}
  {\True Số chấm xuất hiện là số chia hết cho $3$ thì xác xuất  của nó bằng $\dfrac{1}{3}$}
  {Số chấm xuất hiện là số không lớn hơn  $4$ thì xác xuất của nó bằng $\dfrac{5}{6}$}
  {\True Số  chấm xuất hiện là số nguyên tố thì xác xuất của nó bằng $\dfrac{1}{2}$}
  \loigiai{
  Gọi $A\colon$  \lq\lq Mặt một chấm xuất hiện\rq\rq,~ ta có $\mathrm{P}(A)=\dfrac{1}{6}$;\\
  $B\colon$ \lq\lq Mặt hai chấm xuất hiện\rq\rq,~ ta  có $\mathrm{P}(B)=\dfrac{1}{6}$;\\
  $C\colon$ \lq\lq Mặt  ba chấm xuất hiện\rq\rq,~ ta có $\mathrm{P}(C)=\dfrac{1}{6}$;\\
  $D\colon$ \lq\lq Mặt  bốn chấm xuất hiện\rq\rq,~ ta có $\mathrm{P}(D)=\dfrac{1}{6}$;\\
  $E\colon$ \lq\lq Mặt năm  chấm xuất hiện\rq\rq,~ ta có $\mathrm{P}(E)=\dfrac{1}{6}$;\\
  $F\colon$ \lq\lq Mặt sáu chấm xuất hiện\rq\rq,~ ta có $\mathrm{P}(F)=\dfrac{1}{6}$. 
  \begin{itemchoice}
   \itemch Ta có $A\cup C\cup E\colon$ \lq\lq Số chấm xuất hiện là số lẻ\rq\rq, nên suy ra 
   $$\mathrm{P}(A\cup  C\cup E)=\mathrm{P}(A)+\mathrm{P}(C)+\mathrm{P}(E)=\dfrac{1}{2}.$$
   \itemch Ta có $C\cup F\colon$ \lq\lq Số chấm xuất hiện là số chia hết cho $3$\rq\rq,  nên suy ra
   $$\mathrm{P}(C\cup  F)=\mathrm{P}(C)+\mathrm{P}(F)=\dfrac{1}{3}.$$
   \itemch Ta  có $A\cup B\cup C\cup E\colon$ \lq\lq Số chấm xuất hiện là số không lớn hơn $4$\rq\rq, nên suy ra
   $$\mathrm{P}(A\cup B\cup C\cup  D)=\mathrm{P}(A)+\mathrm{P}(B)+\mathrm{P}(C)+\mathrm{P}(D)=\dfrac{2}{3}.$$
   \itemch Ta có $B\cup C\cup E\colon$ \lq\lq Số chấm xuất hiện là số  nguyên tố\rq\rq, nên suy ra
   $$\mathrm{P}(B\cup C\cup  E)=\mathrm{P}(B)+\mathrm{P}(C)+\mathrm{P}(E)=\dfrac{1}{2}.$$
  \end{itemchoice}
  }
\end{ex}

% nguồn: Ok_dungsai_lop11_Lan1
\begin{ex}%[1D9H2-3]%[ma:34beac35]
 Từ tập  hợp $X=\left\{0;2;3;4;5;8\right\}$. Các khẳng định sau đây, khẳng định nào đúng, khẳng định nào sai?
 \choiceTF 
  {Số các số tự nhiên có bốn  chữ số là $4^6$}
  {Số các số tự nhiên có bốn chữ số đôi một khác nhau là  $\mathrm{A}_6^4$}
  {\True Số các số tự nhiên chẵn có  bốn chữ số đôi một khác nhau là $108$}
  {\True Số các số tự  nhiên chẵn có bốn chữ số đôi một khác nhau và nhỏ hơn $2024$ là $50$}
 \loigiai{ 
  \begin{itemchoice}
   \itemch Gọi số có dạng $\overline{abcd}$.\\ 
   Có $5$ cách chọn cho $a$. Có  $6$ cách chọn cho mỗi vị trí $b$, $c$, $d$.\\
   Theo quy tắc nhân, ta có $5\cdot 6\cdot 6\cdot  6=1080$ số thỏa mãn.
   \itemch  Gọi số có dạng $\overline{abcd}$.\\
   Số các số tự nhiên có bốn chữ số đôi một khác nhau  là $\mathrm{A}_6^4-\mathrm{A}_5^3=300$.
   \itemch Gọi  số có dạng $\overline{abcd}$.
   \begin{itemize}
    \item Nếu $d=0$ thì có $\mathrm{A}_5^3=60$  số thỏa mãn.
    \item Nếu $d\in \left\{2;4;8\right\}$ suy ra có $3$ cách chọn cho $d$, có $4$ cách chọn cho $a$ và có $\mathrm{A}_4^2$ cách chọn $b,c$.\\ 
    Theo quy tắc  nhân, ta có $4\cdot 4\cdot \mathrm{A}_4^2=48$ số thỏa mãn.
   \end{itemize}
   Vậy theo quy  tắc cộng, ta có $60+48=108$ số thỏa mãn.
   \itemch Gọi số  có dạng $\overline{abcd}$.
   \begin{itemize}
    \item Nếu  $a=1$ thì $d$ có $4$ cách chọn và $b,c$ có $\mathrm{A}_4^2$ cách chọn.\\
    Theo quy tắc nhân, ta có $1\cdot 4\cdot 4\cdot 3=48$ số thỏa mãn. 
     \item Nếu $a=2$ thì $\heva{&b=0\\ &c=1}\Rightarrow d\in \left\{4;8\right\}.$\\
    Theo quy tắc  nhân, ta có $1\cdot 1\cdot 1\cdot 2=2$ số thỏa mãn.
   \end{itemize}
   Vậy theo quy tắc cộng, ta có  $48+2=50$ số thỏa mãn.
  \end{itemchoice}
 } 
\end{ex}

% nguồn: Ok_dungsai_lop11_Lan1
\begin{ex}%[1D9H2-3]%[ma:94227e8c]
 Một tổ có  $6$ học sinh nam và $9$ học sinh nữ. Xác suất có thể xảy ra khi chọn $6$ học sinh đi lao động. Trong các khẳng định sau, khẳng định nào đúng, khẳng định nào sai?
  \choiceTF
  {\True Xác suất $6$ học sinh được chọn  đều là nam là $\dfrac{1}{5005}$}
  {Xác suất có đúng $2$ học sinh  nam là $\dfrac{3}{143}$}
  {\True Xác suất có cả nam và nữ là $\dfrac{618}{715}$} 
  {Xác  suất có ít nhất $1$ học sinh nữ là $\dfrac{5004}{5005}$}
 \loigiai{ 
  Số cách  chọn ngẫu nhiên $6$ học sinh là $n(\Omega)=\mathrm{C}_{15}^6=5005$.
  \begin{itemchoice}
   \itemch Gọi $A\colon$ \lq\lq $6$ học sinh  được chọn đều là nam\rq\rq~ $\Rightarrow \mathrm{P}(A)=\dfrac{n(A)}{n(\Omega)}=\dfrac{1}{5005}$.
   \itemch Gọi $B\colon$ \lq\lq có đúng  $2$ học sinh nam\rq\rq~ $\Rightarrow \mathrm{P}(B)=\dfrac{n(B)}{n(\Omega)}=\dfrac{\mathrm{C}_6^2\cdot\mathrm{C}_9^4}{\mathrm{C}_{15}^6}=\dfrac{54}{143}$.
   \itemch Gọi $C\colon$ \lq\lq  có cả nam và nữ\rq\rq~ $\Rightarrow \mathrm{P}(C)=\dfrac{n(C)}{n(\Omega)}=\dfrac{\mathrm{C}_6^1\cdot\mathrm{C}_9^5+\mathrm{C}_6^2\cdot\mathrm{C}_9^4+\mathrm{C}_6^3\cdot\mathrm{C}_9^3}{\mathrm{C}_{15}^6}=\dfrac{618}{715}$.
   \itemch Gọi $D\colon$ \lq\lq có ít nhất $1$ học sinh  nữ\rq\rq~ $\Rightarrow \mathrm{P}(D)=\dfrac{n(D)}{n(\Omega)}=1-\dfrac{1}{\mathrm{C}_{15}^6}=\dfrac{5004}{5005}$.
  \end{itemchoice}
 } 
\end{ex}

% nguồn: Ok_dungsai_lop11_Lan1
\begin{ex}%[1D9H2-3]%[ma:bbad31d6]
 Một hộp gồm $6$ viên bi trắng, $4$ viên bi đỏ, $5$  viên bi xanh. Chọn ngẫu nhiên $3$ viên bi. Các khẳng định sau đây, khẳng định nào là đúng, khẳng định nào là sai?
  \choiceTF
  {\True Xác suất  để thu được $3$ viên bi cùng màu là $\dfrac{34}{455}$}
  {Xác suất để thu được  có đúng $2$ viên bi màu xanh là $\dfrac{12}{91}$}
  {\True Xác suất để thu được $3$ viên bi khác màu  là $\dfrac{24}{91}$}
  {Xác  suất để thu được ít nhất $1$ viên bi màu đỏ là $\dfrac{38}{91}$}
  \loigiai{
  Không gian mẫu $n\left(\Omega  \right)=C_{15}^3$.
  \begin{itemchoice}
   \itemch Gọi $A$ là biến cố \lq\lq Lấy $3$  viên bi được chọn có cùng màu\rq\rq.\\
   Gọi $E$  là biến cố \lq\lq Lấy $3$ viên bi cùng màu trắng\rq\rq~ $\Rightarrow \mathrm{P}\left(E \right)=\dfrac{\mathrm{C}_6^3}{\mathrm{C}_{15}^3}=\dfrac{20}{455}$.\\
   Gọi $G$ là biến cố \lq\lq Lấy $3$  viên bi cùng màu đỏ\rq\rq~ $\Rightarrow \mathrm{P}\left(G \right)=\dfrac{\mathrm{C}_4^3}{\mathrm{C}_{15}^3}=\dfrac{4}{455}$.\\
   Gọi $H$ là biến cố \lq\lq Lấy $3$ viên bi cùng  màu xanh\rq\rq~ $\Rightarrow \mathrm{P}\left(H \right)=\dfrac{\mathrm{C}_5^3}{\mathrm{C}_{15}^3}=\dfrac{10}{455}$.\\
   Ta có $E,\,G,\,H$ là các  biến cố xung khắc nên $\mathrm{P}\left(A \right)=\mathrm{P}\left(E \right)+\mathrm{P}\left(G \right)+\mathrm{P}\left(H \right)=\dfrac{34}{455}$.
   \itemch Gọi $B$ là biến cố \lq\lq Lấy  được đúng $2$ viên bi màu xanh\rq\rq.\\
   Gọi  $M$ là biến cố \lq\lq Lấy $3$ viên bi có $2$ bi màu xanh, $1$ viên bi màu đỏ\rq\rq~ $\Rightarrow \mathrm{P}\left(M \right)=\dfrac{\mathrm{C}_5^2\mathrm{C}_4^1}{\mathrm{C}_{15}^3}=\dfrac{40}{455}$.\\
   Gọi $N$  là biến cố \lq\lq Lấy $3$ viên bi có $2$ bi màu xanh, $1$ viên bi màu trắng\rq\rq~ $\Rightarrow \mathrm{P}\left(N \right)=\dfrac{\mathrm{C}_5^2\mathrm{C}_6^1}{\mathrm{C}_{15}^3}=\dfrac{60}{455}$.\\
   Ta có $M,\, N$ là các biến  cố xung khắc nên $\mathrm{P}\left(B \right)=\mathrm{P}\left(M \right)+\mathrm{P}\left(N \right)=\dfrac{40+60}{455}=\dfrac{20}{91}$.
   \itemch Gọi $C$ là biến cố  \lq\lq $3$ viên bi được chọn khác màu\rq\rq~ $\Rightarrow n\left(C \right)=\mathrm{C}_6^1\mathrm{C}_4^1\mathrm{C}_5^1=6\cdot 4\cdot 5=120$.\\
   Xác suất của biến cố $C$ là $\mathrm{P}\left(C \right)=\dfrac{n\left(C \right)}{n\left(\Omega \right)}=\dfrac{6\cdot  4\cdot 5}{\mathrm{C}_{15}^3}=\dfrac{24}{91}$.
    \itemch Gọi $D$ là biến cố \lq\lq Có ít nhất $1$ viên bi màu đỏ\rq\rq.\\
    $\Rightarrow \overline{D}$ là biến cố \lq\lq $3$ viên bi được chọn không có viên bi đỏ nào\rq\rq~ $\Rightarrow \mathrm{P}\left(\overline{D} \right)=\dfrac{\mathrm{C}_{11}^3}{\mathrm{C}_{15}^3}$.\\
   Vậy  $\mathrm{P}\left(D \right)=1-\mathrm{P}\left(\overline{D} \right)=1-\dfrac{\mathrm{C}_{11}^3}{\mathrm{C}_{15}^3}=\dfrac{58}{91}$.
  \end{itemchoice}
  }
\end{ex}

% nguồn: Ok_dungsai_lop11_Lan1
\begin{ex}%[1D9H2-3]%[ma:b0fa7516]
 Một hộp có $18$ viên bi, trong đó có $7$ viên bi mầu đỏ được đánh số từ $1$ đến $7$, $6$ viên bi mầu xanh được đánh số từ $1$ đến $6$ và $5$ viên bi mầu vàng được đánh số từ $1$ đến $5$. Lấy ngẫu nhiên trong hộp ra $3$ viên bi. Các khẳng định sau đây, khẳng định nào đúng, khẳng định nào  sai?
  \choiceTF
  {\True Xác suất để lấy được $3$ viên bi  mầu đỏ là $\dfrac{35}{816}$}
  {Xác suất để lấy được  $3$ viên bi cùng mầu là $\dfrac{67}{816}$}
  {\True Xác suất để lấy được $3$ viên bi đủ  cả ba mầu là $\dfrac{35}{136}$}
  {Xác suất để lấy được $3$ viên  bi khác mầu và khác số là $\dfrac{121}{816}$}
 \loigiai{ 
  Lấy ngẫu nhiên trong hộp ra $3$ viên bi từ $18$ viên bi nên có  $n(\Omega)=\mathrm{C}_{18}^3$.
  \begin{itemchoice}
   \itemch Gọi $A$ là biến cố lấy được  $3$ viên bi màu đỏ. Ta có $n(A)=\mathrm{C}_7^3$.\\
   Nên xác suất để lấy được $3$ viên  bi màu đỏ là $\mathrm{P}(A)=\dfrac{n(A)}{n(\Omega)}=\dfrac{\mathrm{C}_7^3}{\mathrm{C}_{18}^3}=\dfrac{35}{816}$.
   \itemch Gọi $B$ là biến cố lấy được $3$ viên bi cùng mầu. Ta có $n(B)=\mathrm{C}_7^3+\mathrm{C}_6^3+\mathrm{C}_5^3$.\\ 
   Nên xác suất để lấy được $3$ viên bi cùng màu  là $\mathrm{P}(B)=\dfrac{n(B)}{n(\Omega)}=\dfrac{\mathrm{C}_7^3+\mathrm{C}_6^3+\mathrm{C}_5^3}{\mathrm{C}_{18}^3}=\dfrac{65}{816}$.
   \itemch Gọi $C$ là biến cố lấy được $3$ viên  bi đủ cả ba mầu.\\
   Tức mỗi màu có $1$ viên nên ta có xác suất xảy ra biến cố  $C$ là
    $\mathrm{P}(C)=\dfrac{\mathrm{C}_7^1\cdot\mathrm{C}_6^1\cdot\mathrm{C}_5^1}{\mathrm{C}_{18}^3}=\dfrac{210}{816}=\dfrac{35}{136}.$
   \itemch Gọi $D$ là biến cố lấy được $3$ viên bi khác mầu và  khác số.\\
   Khi đó xác suất xảy ra biến cố $D$ là 
   $\mathrm{P}(D)=\dfrac{5\cdot 5\cdot  5}{\mathrm{C}_{18}^3}=\dfrac{125}{816}.$
  \end{itemchoice}
 } 
\end{ex}

% nguồn: Ok_dungsai_lop11_Lan1
\begin{ex}%[1D9H2-3]%[ma:7febe37e]
 Hai bạn Y và Q mỗi bạn chọn ngẫu nhiên  một số nguyên dương có một chữ số. Các khẳng định sau đây, khẳng định nào là đúng, khẳng định nào là sai?
  \choiceTF
  {\True  Xác suất để bạn Y chọn được số là số chẵn là $\dfrac{4}{9}$}
   {Xác suất để bạn Q chọn được số là số lẻ và chia hết cho $3$ là $\dfrac{1}{3}$}
  {Xác suất để hai bạn Y và Q chọn được hai số  giống nhau là $\dfrac{4}{9}$}
  {\True Xác suất để hai bạn Y  và Q chọn được hai số có tích là một số chẵn là $\dfrac{56}{81}$}
  \loigiai{
  Số các số nguyên dương có một chữ số là $9$ số, trong đó có $4$ số nguyên dương chẵn, $5$ số nguyên dương lẻ, $2$ số lẻ và  chia hết cho 3.\\
  \begin{itemchoice}
   \itemch Xác suất để bạn Y chọn được số là số  chẵn bằng $\dfrac{4}{9}$.
   \itemch Xác suất để bạn Q chọn được số là số lẻ và  chia hết cho $3$ bằng $\dfrac{2}{9}$.
   \itemch Bạn Y chọn chọn ngẫu nhiên  một số nguyên dương có một chữ số có $9$ cách chọn.\\
   Bạn Q chọn ngẫu nhiên một số nguyên dương giống bạn Y đã chọn có  $1$ cách chọn.\\
   Gọi $A$  là biến cố \lq\lq Hai bạn Y và Q chọn được hai số giống nhau\rq\rq.\\
   Số khả năng thuận  lợi của biến cố $A$ là $n(A)=9\cdot 1=9$.\\
   Số phần tử của  không gian mẫu là $n\left(\Omega \right)=9\cdot 9=81$.\\
   Vậy xác suất để hai bạn Y và Q chọn được hai số giống nhau $\mathrm{P}\left(A \right)=\dfrac{n\left(A \right)}{n\left(\Omega  \right)}=\dfrac{9}{81}=\dfrac{1}{9}$.
   \itemch  Gọi $B$ là biến cố \lq\lq Y và Q chọn được hai số có tích là một số chẵn\rq\rq.\\
    Hai số có tích là một số chẵn khi và chỉ khi ít nhất một trong hai số là số chẵn.\\
   Ta xét các biến  cố sau đây
   \begin{itemize}
    \item $B_1\colon$ \rq\rq Y và Q chọn được hai số cùng là số chẵn\rq\rq~ $\Rightarrow n\left(B_1 \right)=4\cdot  4=16 \Rightarrow \mathrm{P}\left(B_1 \right)=\dfrac{n\left(B_1 \right)}{n\left(\Omega \right)}=\dfrac{16}{81}$.
    \item $B_2\colon$ \rq\rq Y chọn được số chẵn và Q chọn được số lẻ\rq\rq~ $\Rightarrow n\left(B_1 \right)=4\cdot 5=20  \Rightarrow \mathrm{P}\left(B_2 \right)=\dfrac{n\left(B_2 \right)}{n\left(\Omega \right)}=\dfrac{20}{81}$.
    \item $B_3\colon$ \rq\rq Y chọn được số lẻ và Q chọn được số chẵn\rq\rq~ $\Rightarrow n\left(B_3 \right)=5\cdot  4=20 \Rightarrow \mathrm{P}\left(B_3 \right)=\dfrac{n\left(B_3 \right)}{n\left(\Omega \right)}=\dfrac{20}{81}$.
   \end{itemize}
   Suy ra $B_1,\,B_2,\,B_3$ đôi một xung khắc và $B=B_1\cup B_2\cup B_3$.\\ 
   Xác suất để hai bạn Y và Q chọn được hai số có tích là một  số chẵn là
    $$\mathrm{P}(B)=\mathrm{P}(B_1)+\mathrm{P}(B_2)+\mathrm{P}(B_3)=\dfrac{16}{81}+\dfrac{20}{81}+\dfrac{20}{81}=\dfrac{56}{81}.$$
  \end{itemchoice}
 } 
\end{ex}

% nguồn: Ok_dungsai_lop11_Lan1
\begin{ex}%[1D9H2-3]%[ma:1ae71826]
  Có $40$ tấm thẻ được đánh số từ $1$ đến $40$, rút ngẫu nhiên ba tấm thẻ. Trong các khẳng định sau đây, khẳng định nào đúng, khẳng định nào sai?
  \choiceTF
  {Biến cố $A\colon$ \lq\lq rút được ba thẻ được đánh số lẻ\rq\rq và $B\colon$ \lq\lq rút được ba thẻ trong đó có đúng một thẻ được đánh sổ  lẻ\rq\rq~ là hai biến cô không xung khắc}
  {\True Xác suất của  biến cố $C\colon$ \lq\lq rút được ba thẻ mà tổng ba số được đánh trên đó là một số lẻ\rq\rq~ là $\dfrac{1}{2}$}
  {Xác suất của  biến cố $D\colon$ \lq\lq rút được ba thẻ mà tích ba số trên đó là một số chẵn\rq\rq~ là $\dfrac{3}{26}$}
  {\True Xác suất của biến cố $E\colon$ \lq\lq rút được ba thẻ mà tổng ba số được đánh  trên đó là một số chia hết cho $3$\rq\rq~ là $\dfrac{127}{380}$}
 \loigiai{ 
  \begin{itemchoice}
   \itemch Ta có hai biến cố $A$ và $B$ là hai biến cố xung khắc vì nếu biến cố $A$ xảy ra thì  biến cố $B$ sẽ không xảy ra và ngược lại.
   \itemch Ta có tổng ba số là một số lẻ khi và chỉ khi cả ba số đó đều là số lẻ hoặc trong ba số chỉ có một số lẻ.  Do đó $C=A\cup B$.\\
   Ta có $n_{\Omega}=\mathrm{C}_{40}^3$  và $n(A)=\mathrm{C}_{20}^3,\,n(B)=\mathrm{C}_{20}^1\cdot \mathrm{C}_{20}^2$.\\
   Suy ra $\mathrm{P}(A)=\dfrac{n(A)}{n(\Omega)}=\dfrac{\mathrm{C}_{20}^3}{\mathrm{C}_{40}^3}=\dfrac{3}{26}$ và $\mathrm{P}(B)=\dfrac{n(B)}{n(\Omega)}=\dfrac{\mathrm{C}_{20}^1\cdot  \mathrm{C}_{20}^2}{\mathrm{C}_{40}^3}=\dfrac{5}{13}$.\\
   Vì $A$ và $B$ là hai biến cố xung khắc  nên $\mathrm{P}(C)=\mathrm{P}(A\cup B)=\mathrm{P}(A)+\mathrm{P}(B)=\dfrac{3}{26}+\dfrac{5}{13}=\dfrac{1}{2}$.
   \itemch  Ta có tích ba số là một số lẻ khi và chỉ khi cả ba số đều là số lẻ.\\
   Do đó hai biến cố $A$  và $D$ là hai biến cố đối (và cũng là xung khắc nhau) cho nên
    $$\mathrm{P}(D)=1-\mathrm{P}(A)=1-\dfrac{3}{26}=\dfrac{23}{26}.$$
   \itemch Gọi các tập hợp  $E_0=\left\{x=3k\left|k=\overline{1,13} \right.\right\}$, $E_1=\left\{x=3k+1\left|k=\overline{0,13} \right.\right\}$, $E_2=\left\{x=3k+2\left|k=\overline{0,12} \right.\right\}$.\\
   Gọi $M$ là biến cố rút được ba thẻ mà  mỗi số trên đó thuộc mỗi tập hợp $E_0,\,E_1,\,E_2$.\\
   $N$ là biến cố rút được ba thẻ mà trong đó cả ba số trên đó đều thuộc  cùng một tập hợp.\\
   Vì $M,\,N$ xung khắc nhau và $E=M\cup N$ nên 
   $$\mathrm{P}(E)=\mathrm{P}(M\cup N)=\mathrm{P}(M)+\mathrm{P}(N)=\dfrac{\mathrm{C}_{13}^1\cdot  \mathrm{C}_{14}^1\cdot \mathrm{C}_{13}^1}{\mathrm{C}_{40}^3}+\dfrac{\mathrm{C}_{13}^3+\mathrm{C}_{14}^3+\mathrm{C}_{13}^3}{\mathrm{C}_{40}^3}=\dfrac{127}{380}.$$
  \end{itemchoice}
 } 
\end{ex}

% nguồn: Ok_dungsai_lop11_Lan1
\begin{ex}%[1D9V2-3]%[ma:d5e4c683]
 Hộp I đựng $4$ viên bi màu đỏ và $6$ viên bi màu xanh, hộp II đựng $2$ viên bi màu đỏ và $8$ viên bi màu xanh. Gieo một con súc sắc, nếu được mặt $6$ chấm thì lấy $1$ viên bi từ hộp I, nếu được mặt khác thì lấy $1$ viên từ hộp II. Gọi $A$  là biến cố \lq\lq Con súc sắc xuất hiện mặt $6$\rq\rq, $B_i$ là biến cố \lq\lq Lấy được viên bi màu đỏ từ hộp $i$\rq\rq, $\left(i=1,2\right)$.
 \choiceTF 
  {\True Xác suất để  gieo con súc sắc xuất hiện mặt $6$ chấm là $\mathrm{P}\left(A \right)=\dfrac{1}{6}$}
  {\True Xác suất  lấy được bi màu đỏ từ hộp I là $\mathrm{P}\left(B_1 \right)=\dfrac{2}{5}$}
  {\True Xác suất lấy đc bi  từ hộp II là $\mathrm{P}\left(B_2 \right)=\dfrac{1}{5}$}
   {Xác suất lấy được bi màu đỏ là $\dfrac{3}{5}$}
  \loigiai{
  \begin{itemchoice}
   \itemch Ta có xác suất để gieo con súc sắc xuất hiện mặt $6$  chấm là $\mathrm{P}\left(A \right)=\dfrac{1}{6}$.
    \itemch Xác suất để lấy từ hộp I được bi màu đỏ là $\mathrm{P}\left(B_1 \right)=\dfrac{4}{10}=\dfrac{2}{5}$.
   \itemch Xác  suất để lấy từ hộp II được bi màu đỏ là $\mathrm{P}\left(B_2 \right)=\dfrac{2}{10}=\dfrac{1}{5}$.
   \itemch Gọi $C$ là biến cố \lq\lq Lấy được viên bi màu đỏ\rq\rq~ suy ra $C=AB_1 \cup \overline{A}B_2$.\\ 
   Xác suất để gieo con súc sắc không xuất hiện mặt $6$ chấm là  $\mathrm{P}\left(\overline{A} \right)=\dfrac{5}{6}$.\\
   Vậy $\mathrm{P}(C)=\mathrm{P}\left(A  \right)\cdot \mathrm{P}\left(B_1 \right)+\mathrm{P}\left(\overline{A} \right)\cdot\mathrm{P}\left(B_2 \right)=\dfrac{1}{6}\cdot\dfrac{2}{5}+\dfrac{5}{6}\cdot\dfrac{1}{5}=\dfrac{7}{30}$.
  \end{itemchoice}
 } 
\end{ex}

% nguồn: Ok_dungsai_lop11_Lan1
\begin{ex}%[1D9V2-3]%[ma:107fa3b7]
 Một hộp đựng $10$ bi trong đó có $4$ bi  xanh và $6$ bi đỏ. Lấy ra $1$ bi và không hoàn lại, tiếp tục lấy ra $1$ bi nữa. Xét các biến cố\\
  \indent \indent $A\colon$ \lq\lq Bi lấy được lần $1$ màu xanh\rq\rq;\\
 \indent \indent $B\colon$ \lq\lq Bi lấy  được lần $2$ màu đỏ\rq\rq.\\
 Trong các  khẳng định sau đây, khẳng định nào đúng, khẳng định nào sai?
 \choiceTF 
  {$A,B$ là hai biến  cố xung khắc}
  {$\mathrm{P}\left(AB  \right)=\dfrac{6}{25}$}
  {$\mathrm{P}\left(AB \right)=\mathrm{P}\left(A  \right)\cdot \mathrm{P}\left(B \right)+\dfrac{2}{25}$}
  {\True Xác suất để $2$ bi lấy được cùng màu  là $\dfrac{7}{15}$}
  \loigiai{
  \begin{itemchoice}
   \itemch Do $A,B$ có  thể cùng đồng thời xảy ra nên $A,B$ không phải hai biến cố xung khắc.
   \itemch Ta  có $n\left(\Omega \right)=9\cdot 10=90$.\\
   $AB\colon$ \lq\lq Bi lần $1$ màu xanh và bi lần $2$ màu đỏ\rq\rq $\Rightarrow n\left(AB  \right)=4\cdot 6=24$.\\
   Do đó  $\mathrm{P}\left(AB \right)=\dfrac{24}{90}=\dfrac{4}{15}$.
   \itemch Ta có $\mathrm{P}\left(A  \right)=\dfrac{2}{5}$.
   \begin{itemize}
    \item Lần $1$ lấy được màu xanh, lần $2$ lấy được màu  đỏ có $4\cdot 6=24$ (cách).
    \item Lần $1$ lấy được màu đỏ, lần $2$ lấy  được màu đỏ có $6\cdot 5=30$ (cách).
   \end{itemize}
   Suy ra $\mathrm{P}\left(B \right)=\dfrac{24+30}{90}=\dfrac{3}{5}$.\\ 
   Do đó $\mathrm{P}\left(A \right)\cdot \mathrm{P}\left(B \right)=\dfrac{6}{25} \Rightarrow \mathrm{P}\left(AB \right) \ne \mathrm{P}\left(A \right)\cdot  \mathrm{P}\left(B \right)+\dfrac{2}{25}$.
   \itemch $C\colon$ \lq\lq  Hai bi lấy được cùng màu\rq\rq~ suy ra $C=A\overline{B}\cup \overline{A}B$.\\
   Rõ ràng $A\overline{B}$  và $\overline{A}B$ là xung khắc nhau.\\
   Ta  có $\mathrm{P}\left(A\overline{B} \right)=\dfrac{4\cdot 3}{90}=\dfrac{2}{15};\,\mathrm{P}\left(\overline{A}B \right)=\dfrac{6\cdot 5}{90}=\dfrac{1}{3}$.\\
   Vậy  $\mathrm{P}\left(C \right)=\mathrm{P}\left(A\overline{B} \right)+\mathrm{P}\left(\overline{A}B \right)=\dfrac{2}{15}+\dfrac{1}{3}=\dfrac{7}{15}$.
  \end{itemchoice}
 } 
\end{ex}

% nguồn: Ok_dungsai_lop11_Lan1
\begin{ex}%[1D9V2-3]%[ma:5be49be7]
 Trong một bài thi trắc nghiệm khách quan gồm $10$ câu. Mỗi câu có $4$ phương án trả lời, trong đó chỉ có một phương án đúng. Một học sinh chuẩn bị bài không tốt nên làm bài bằng cách: Với mỗi câu,  chọn ngẫu nhiên một phương án trả lời. Xét tính đúng sai trong các khẳng định sau.
 \choiceTF 
  {\True Xác suất để học sinh đó trả lời đúng  mỗi câu là $0{,}25$}
  {Xác suất để học sinh đó trả lời sai cả $10$ là  $10 \cdot 0{,}75$}
  {\True Xác suất  để học sinh đó trả lời sai mỗi câu là $0{,}75$}
  {\True Xác suất để học sinh đó trả lời đúng $1$ câu là $10 \cdot 0{,}25$} 
  \loigiai{
 Mỗi câu, xác suất để học sinh  chọn sai là $\dfrac{3}{4} = 0{,}75$.\\
  Mỗi câu, xác suất để học sinh chọn đúng là $\dfrac{1}{4} = 0{,}25$.\\
 Khi đó, xác suất để học sinh chọn sai cả 10  câu là
  \[\underbrace{\dfrac{3}{4}\cdot\dfrac{3}{4}\cdots\cdot\dfrac{3}{4}}_{10}=\left(\dfrac{3}{4}\right)^{10}=(0{,}75)^{10}.\]
 Xác suất để học sinh chọn đúng $1$ câu  là
 \[\mathrm{C}_{10}^1 \cdot  \dfrac{1}{4} = 10 \cdot0{,}25. \]
 } 
\end{ex}

% nguồn: Ok_dungsai_lop11_Lan1
\begin{ex}%[1D9V2-3]%[ma:f8f783fc]
 Một bài trắc nghiệm có $10$ câu hỏi, mỗi câu hỏi có $4$ phương án lựa chọn trong đó có $1$ đáp án đúng. Giả sử mỗi câu trả lời đúng được cộng $5$ điểm và mỗi câu trả lời sai bị trừ đi $2$ điểm. Một học sinh không học bài nên chọn hú họa một đáp án.  Trong các khẳng định sau, khẳng định nào đúng, khẳng định nào sai?
 \choiceTF 
  {\True Số câu học sinh đó trả lời đúng để được $1$ điểm  là $3$}
  {Xác suất để học sinh trả lời được $1$ điểm là $\left(\dfrac{1}{4} \right)^3\cdot\left(\dfrac{3}{4} \right)^7$} 
  {\True Xác suất để học  sinh không trả lời được câu nào là $\left(\dfrac{3}{4} \right)^{10}$}
  {\True Xác suất để học sinh đó nhận điểm không lớn hơn $1$ là $\approx 0{,}7759$} 
 \loigiai{ 
  \begin{itemchoice}
   \itemch Gọi  $x$ là số câu học sinh đó trả lời đúng.\\
   Theo đề bài ta có học sinh đó  nhận điểm $1$, suy ra $5x-2\left(10-x \right)=1 \Leftrightarrow 7x=21 \Leftrightarrow x=3$.
   \itemch Xác suất trả lời  đúng của học sinh trong một câu là $\dfrac{1}{4}$.\\
   Xác suất trả lời sai của học  sinh trong một câu là $\dfrac{3}{4}$.\\
   Học sinh trả  lời đúng $3$ câu là $\mathrm{P}_3=\mathrm{C}_{10}^3\cdot\left(\dfrac{1}{4} \right)^3\cdot\left(\dfrac{3}{4} \right)^7$.
   \itemch Học sinh trả lời không đúng  câu nào $\mathrm{P}_0=\left(\dfrac{3}{4} \right)^{10}$.
   \itemch Theo câu a, ta có học sinh đó nhận điểm không lớn hơn $1$ thì học sinh này  trả lời đúng nhiều nhất là $3$ câu.\\
    Ta có
   \begin{itemize}
    \item Học sinh trả  lời không đúng câu nào $\mathrm{P}_0=\left(\dfrac{3}{4} \right)^{10}$.
    \item Học sinh trả lời đúng  $1$ câu $\mathrm{P}_1=\mathrm{C}_{10}^1\cdot\left(\dfrac{1}{4} \right)^1\cdot\left(\dfrac{3}{4} \right)^9$.
    \item Học sinh  trả lời đúng $2$ câu $\mathrm{P}_2=\mathrm{C}_{10}^2\cdot\left(\dfrac{1}{4} \right)^2\cdot\left(\dfrac{3}{4} \right)^8$.
    \item Học sinh trả lời đúng $3$  câu $\mathrm{P}_3=\mathrm{C}_{10}^3\cdot\left(\dfrac{1}{4} \right)^3\cdot\left(\dfrac{3}{4} \right)^7$.
   \end{itemize}
   Vậy xác suất cần tìm là  $\mathrm{P}=\mathrm{P}_0+\mathrm{P}_1+\mathrm{P}_2+\mathrm{P}_3 \approx 0,7759$.
  \end{itemchoice}
 } 
\end{ex}

% nguồn: Ok_dungsai_lop11_Lan1
\begin{ex}%[1D9V2-3]%[ma:440d671d]
 Một hộp có $5$ quả bóng xanh, $6$ quả bóng đỏ và $4$ quả bóng vàng có kích thước và khối lượng như nhau. Lấy ngẫu  nhiên từ hộp $4$ quả bóng.
  \choiceTF
  {\True Xác suất của biến cố \lq\lq Cả $4$ quả bóng lấy ra cùng màu\rq\rq\, là $\dfrac{1}{65}$} 
  {\True  Xác suất của biến cố \lq\lq Trong $4$ quả bóng lấy ra có đủ $3$ màu\rq\rq\, là $\dfrac{48}{91}$}
  {\True Số phần tử không gian mẫu là  $1365$.}
  {Xác suất của  biến cố \lq\lq Cả $4$ quả bóng lấy ra cùng màu\rq\rq\, lớn hơn $0{,}5$}
 \loigiai{ 
  Số phần tử không gian mẫu là $n(\Omega)  = \mathrm{C}_{15}^4 = 1365$.
  \begin{itemize}
   \item Xét biến cố  $A\colon$ \lq\lq Cả $4$ quả bóng lấy ra cùng màu \rq\rq\\
   Suy  ra $n(A) = \mathrm{C}_{5}^4 + \mathrm{C}_{6}^4 + \mathrm{C}_{4}^4 = 21$.\\
   Vậy $P(A) = \dfrac{n(A)}{n(\Omega)} = \dfrac{21}{1365}=  \dfrac{1}{65}$.
   \item Xét biến cố $B\colon$ \lq\lq  Trong $4$ quả bóng lấy ra có đủ $3$ màu \rq\rq\\
   Suy ra $n(B) = \mathrm{C}_{5}^2 \cdot \mathrm{C}_{6}^1 \cdot \mathrm{C}_{4}^1 +  \mathrm{C}_{5}^1 \cdot \mathrm{C}_{6}^2 \cdot \mathrm{C}_{4}^1 + \mathrm{C}_{5}^1 \cdot \mathrm{C}_{6}^1 \cdot \mathrm{C}_{4}^2 = 720$.\\
   Vậy  $P(A) = \dfrac{n(A)}{n(\Omega)} = \dfrac{720}{1365} = \dfrac{48}{91}$.
  \end{itemize}
  \begin{itemchoice}
   \itemch  Đúng
   \itemch  Đúng
    \itemch Đúng
   \itemch  Sai
  \end{itemchoice}
 } 
\end{ex}

% nguồn: Ok_dungsai_lop11_Lan1
\begin{ex}%[1D9V2-3]%[ma:7eebd3dc]
 Trên một bảng quảng cáo, người ta mắc hai hệ thống bóng đèn. Hệ thống I gồm 2 bóng mắc nối  tiếp, hệ thống II gồm 2 bóng mắc song song. Khả năng bị hỏng của mỗi bóng đèn sau 6 giờ thắp sáng liên tục là 0{,}15. Biết tình trạng của mỗi bóng đèn là độc lập. Khi đó xác suất để
  \choiceTF
  {\True Hệ thống II bị hỏng  (không sáng) bằng $0{,}0225$}
   {\True Hệ thống II hoạt động bình thường bằng $0{,}9775$}
   {Hệ thống I bị hỏng (không sáng) bằng $0{,}5775$}
  {Cả hai hệ thống bị hỏng (không sáng) bằng  $\approx 0{,}02624$}
  \loigiai{
  \begin{itemchoice}
   \itemch Đúng. Nhận xét: Hệ thống II gồm  2 bóng được mắc song song nên nó chỉ hỏng khi cả hai bóng đều hỏng.\\
   Gọi $B$ là biến cố: \lq\lq Hệ thống II bị hỏng\rq\rq, ta có:  $\mathrm{P}(B)=0{,}15\cdot 0{,}15=0{,}0225$.
   \itemch Đúng. Xác suất để hệ thống II hoạt  động bình thường: $\mathrm{P}(\overline{B})=1-0{,}0225=0{,}9775$.
   \itemch Sai. Nhận xét: Hệ thống I chỉ hoạt  động bình thường khi cả hai bóng bình thường.\\
   Gọi $A$ là biến cố: "Hệ thống I bị hỏng". Khi đó xác suất để hệ thống I hoạt  động bình thường là $\mathrm{P}(\overline{A})=0{,}85\cdot 0{,}85=0{,}7225$.\\
   Suy ra  $\mathrm{P}(A)=1-0{,}7225=0{,}2775$.
   \itemch Sai. Xác suất để cả hai hệ thống I, II đều  bị hỏng là\\
   $\mathrm{P}(AB)=\mathrm{P}(A)\cdot\mathrm{P}(B)=0{,}2775\cdot 0{,}0225=\dfrac{999}{160000}\approx 0{,}00624$. 
  \end{itemchoice}
 } 
\end{ex}


%%%%%%%%%%%%%%%% TRẢ LỜI NGẮN %%%%%%%%%%%%%%%%

% nguồn: Ok_TraLoiNgan_Lop11_Lan1
\begin{ex}%[1D9H2-3]%[ma:2c3542ca]
	Một chiếc máy có hai động cơ $I$ và $II$ chạy độc lập nhau. Xác suất để động cơ $I$ và $II$ chạy  tốt lần lượt là $0{,}8$ và $0{,}7.$ Tính xác suất để ít nhất một động cơ chạy tốt.
	 \shortans{0,94}
	\loigiai{ 
		Gọi $A, B$  là xác xuất động cơ $I, II$ chạy tốt.\\
		Xác xuất ít nhất một động cơ  chạy tốt là
		\begin{align*}
			P(A\cup B)=P(A)+P(B)-P(A\cap B)=P(A)+P(B)-P(A)P(B)=0{,}8+0{,}7-0{,}8\cdot0{,}7=0{,}94. 
		\end{align*}
	 }
\end{ex}

% nguồn: Ok_TraLoiNgan_Lop11_Lan1
\begin{ex}%[1D9V2-3]%[ma:a02908ce]
	Một bệnh truyền nhiễm có xác suất truyền bệnh là $0{,}7$ nếu tiếp xúc với người bệnh mà không đeo khẩu trang là $0{,}2$ nếu tiếp xúc với người bệnh mà đeo khẩu trang. Xác suất anh An ít nhất một lần bị lây bệnh từ người bệnh mà anh tiếp  xúc người bệnh $5$ lần đều không mang khẩu trang là bao nhiêu? (kết quả làm tròn đến chữ số thập phân thứ ba)
	\shortans{1} 
	\loigiai{Gọi biến cố $A$ : \lq\lq Anh An ít nhất một lần  bị lây bệnh khi tiếp xúc người bệnh cả $5$ lần đều không mang khẩu trang \rq\rq.
		 Biến cố $\bar{A}$ : \lq\lq Anh An không bị lây bệnh khi tiếp xúc người bệnh cả $5$ lần đều không mang khẩu trang\rq\rq. \\
		Xác suất nhiễm bệnh nếu tiếp xúc với người  bệnh mà không đeo khẩu trang là $0{,}7$ .
		 Xác suất không bị nhiễm bệnh nếu tiếp xúc với người bệnh mà không đeo khẩu trang là $1-0{,}7=0{,}3$. \\
	 Từ đó $\mathrm{P}(\bar{A})=(0{,}3)^5 \Rightarrow \mathrm{P}(A)=1-\mathrm{P}(\bar{A})=1-(0{,}3)^5=0{,}998$.}
\end{ex}

% nguồn: Ok_TraLoiNgan_Lop11_Lan1
\begin{ex}%[1D9V2-3]%[ma:d1aa71b5]
	Một hộp chứa $5$ viên bi đỏ, $6$ viên bi xanh và $7$ viên bi vàng (các viên bi kích thước như nhau). Lấy ngẫu nhiên $4$ viên bi từ hộp. Tính xác suất để trong $4$ viên bi lấy được có nhiều nhất 2 viên  bi đỏ (làm tròn kết quả đến hàng phần trăm).
	 \shortans{0,39}
	\loigiai{ 
		 Lấy ngẫu nhiên $4$ viên bi từ hộp. Số kết quả có thể xảy ra là $n({\Omega}_1)=\mathrm{C}_{18}^4$.\\
		Gọi $A$ là biến cố: \lq\lq $4$ viên bi lấy được trong hộp có nhiều nhất 2 viên bi đỏ\rq\rq .\begin{itemize}
		\item[TH1:] Không lấy được viên bi đỏ nào  trong $4$ bi có $\mathrm{C}_5^{0}\cdot\mathrm{C}_{13}^4$ cách.
		\item[TH2:] Lấy được $1$  bi đỏ trong $4$ bi có $\mathrm{C}_5^1\cdot\mathrm{C}_{13}^3$ cách.
		 \item[TH3:] Lấy được $2$ bi đỏ trong $4$ bi có $\mathrm{C}_5^2\cdot\mathrm{C}_{13}^2$ cách.
	\end{itemize}
	 Số kết quả thuận lợi cho $A$ là
	$n(A)=\mathrm{C}_5^0\cdot\mathrm{C}_{13}^4+\mathrm{C}_5^1\cdot\mathrm{C}_{13}^3+\mathrm{C}_5^2\cdot\mathrm{C}_{13}^2$.\\ 
	Khi đó, xác suất để trong $4$ viên bi lấy  được trong hộp có nhiều nhất 2 viên bi đỏ là\\
 $\mathrm{P}(A)=\dfrac{n(A)}{n({\Omega}_1)}=\dfrac{\mathrm{C}_5^0\cdot\mathrm{C}_{13}^4+\mathrm{C}_5^1\cdot\mathrm{C}_{13}^3+\mathrm{C}_5^2\cdot\mathrm{C}_{13}^2}{\mathrm{C}_{18}^4}=\dfrac{481}{1224}=0{,}39$.}
\end{ex}

% nguồn: Ok_TraLoiNgan_Lop11_Lan1
\begin{ex}%[1D9V2-3]%[ma:6c3a9a2e]
	Cường, Trọng cùng $6$ bạn nữ xếp ngẫu nhiên thành một hàng ngang để chụp ảnh. Tính xác suất của biến cố \lq\lq ít nhất một trong  hai bạn Cường và Trọng đứng đầu hàng\rq\rq\, (làm tròn kết quả đến hàng phần trăm).
	\shortans{0,46} 
	\loigiai{ 
		Số phần tử không gian mẫu là $n(\Omega) =  8! = 40320$.\\
		Gọi  $A$ là biến cố \lq\lq ít nhất một trong hai bạn Cường và Trọng đứng đầu hàng\rq\rq.\\
		Suy ra  $\overline{A}$ là \lq\lq Hai bạn Cường và Trọng đều không đứng đầu hàng\rq\rq.\\
		Ta có $n\left(\overline{A}\right) = \mathrm{C}_6^2 \cdot 2! \cdot 6! = 21600$.\\ 
		Vậy xác suất  của biến cố $A$ là $ \mathrm{P}(A) = 1 - \mathrm{P}\left(\overline{A}\right) = 1 - \dfrac{21600}{40320} = \dfrac{13}{28}\approx 0{,}46$.
	 }
\end{ex}

% nguồn: Ok_TraLoiNgan_Lop11_Lan1
\begin{ex}%[1D9V2-3]%[ma:aa4268eb]
	Một chiếc hộp có chín thẻ đánh số thứ tự từ $1$ đến $9$. Rút ngẫu nhiên $2$ thẻ rồi nhân hai số ghi trên thẻ lại với  nhau. Tính xác suất để kết quả nhân được là một số chẵn (kết quả làm tròn đến hàng phần trăm).
	\shortans{0,72} 
	 \loigiai{
		Trường hợp 1: hai số rút ra  đều là số chẵn: $P_1=\dfrac{\mathrm{C}_4^2}{\mathrm{C}_9^2}=\dfrac{1}{6}$.\\
		Trường hợp 2:  hai số rút ra có một số lẻ, một số chẵn: $P_2=\dfrac{\mathrm{C}_4^1\cdot\mathrm{C}_5^1}{\mathrm{C}_9^2}=\dfrac{5}{9}$.\\
	Vậy xác suất để kết quả nhân được là một số  chẵn là $P=P_1+P_2=\dfrac{1}{6}+\dfrac{5}{9}=\dfrac{13}{18}$.}
\end{ex}
